\capitulo{1}{Introducción}

La contratación pública~\cite{inroduccion:contr_publica} constituye uno de los pilares fundamentales del sector público. A través de este proceso, las administraciones adquieren bienes, servicios y obras necesarias para el desarrollo de sus competencias, garantizando los principios de transparencia, libre concurrencia, igualdad de trato y eficiencia en el uso de los recursos públicos. En España, este procedimiento se materializa en la publicación de licitaciones que incluyen pliegos administrativos y técnicos, resoluciones, anexos y diversa documentación normativa que regula cada contrato.

Sin embargo, aunque la información es pública y accesible, su volumen, complejidad técnica y extensión documental dificultan enormemente su creación, consulta y análisis. Los pliegos de contratación pueden superar fácilmente decenas o incluso cientos de páginas, contienen lenguaje jurídico-administrativo especializado y presentan estructuras heterogéneas. Para empresas licitadoras, profesionales del sector o incluso para ciudadanos interesados, localizar información concreta puede convertirse en una tarea lenta y costosa.

En este contexto surge la necesidad de herramientas inteligentes que permitan consultar grandes volúmenes de documentación pública de forma eficiente, precisa y contextualizada. Los modelos de lenguaje de gran tamaño (\eng{LLM}) han demostrado un gran potencial para la generación de texto y la comprensión del lenguaje natural. Sin embargo, presentan limitaciones importantes, ya que su conocimiento es estático (fecha de corte), pueden generar respuestas incorrectas con aparente seguridad (alucinaciones) y no tienen acceso directo a información específica o actualizada, como documentos concretos de licitación.

Para solventar estas limitaciones, en este Trabajo Fin de Grado se ha desarrollado \eng{PythIA}. Un proyecto que propone el diseño e implementación de un sistema basado en la técnica \eng{Retrieval-Augmented Generation} (\eng{RAG}), cuyo objetivo es combinar la capacidad generativa de un \eng{LLM} con un sistema de recuperación semántica de información sobre una colección documental formada por licitaciones públicas. El sistema permite indexar automáticamente los documentos, fragmentarlos en unidades manejables (\eng{chunks}), representarlos mediante vectores lineales (\eng{embeddings}) y almacenarlos en una base de datos vectorial especializada (\eng{Qdrant}). Ante una consulta del usuario, el sistema recupera los fragmentos más relevantes y los incorpora al \eng{prompt} del modelo generativo, mejorando la precisión, actualidad y fiabilidad de la respuesta.

El proyecto no se limita al diseño conceptual del sistema \eng{RAG}, sino que aborda su implementación completa en un entorno realista y desplegable. Para ello, se ha desarrollado una aplicación \eng{web} con \eng{Flask} que permite la interacción del usuario con el sistema de forma transparente, se integra una base de datos relacional para la gestión de usuarios y consultas, y se utiliza una base de datos vectorial (\eng{Qdrant}) para la recuperación semántica. Además, se emplea \eng{Docker} para garantizar la reproducibilidad del entorno y facilitar el despliegue de la aplicación en producción.

Este documento es la memoria principal del proyecto. En él se describe el contexto en el que surge el proyecto, y los objetivos que se persiguen conseguir. Se incluye una explicación de los fundamentos teóricos necesarios para comprender el proyecto, y de las herramientas empleadas para su desarrollo indicando el valor que aportan al proyecto frente a otras alternativas del mercado. También, se exponen los aspectos más relevantes que tuvieron lugar durante el desarrollo explicando sus dificultades y como se solventaron. Para poner en valor real el proyecto es importante conocer el estado del arte en el campo de los modelos de lenguaje alimentados con \eng{RAG}, explicando y comparando el proyecto con otros similares ya existentes. Finalmente, se detallan las conclusiones obtenidas tras finalizar el proyecto y se proponen vías de mejoras para futuros desarrollos.

Junto con la memoria se entrega los siguientes materiales adjuntos:

\begin{itemize}
\tightlist
	\item Documentación técnica: detallada en los anexos. En ella se describe la planificación temporal del proyecto, explicando la planificación de cada \eng{sprint} y comparándola con lo que se realizó en realidad. Se realiza un estudio sobre la viabilidad del proyecto teniendo en cuenta tanto si es económicamente rentable como su viabilidad legal. Se especifican tanto los requisitos, funcionales y no funcionales explicando de manera detalla los casos de uso del proyecto, como el diseño de las distintas partes de la aplicación (arquitectura, datos, procedimientos e interfaces). En este documento, también, se incluyen los manuales del programador y del usuario. Finalmente, se realiza un estudio sobre la sostenibilidad del proyecto.
	\item \eng{PythIA}: página \eng{web} del proyecto. Accesible a través de la dirección de dominio:~\url{https://pythia.es/}
	\item Repositorio~\myhref{https://github.com/lbr1014/TFG_RAG.git}{\eng{GitHub}} del proyecto. Donde se encuentra el código, la documentación y el desarrollo del proyecto. Pudiendo verse en él la planificación, ya que cada \eng{milestone} se corresponde con un \eng{sprint}.
		\item Progreso del proyecto: ha sido gestionado utilizando~\myhref{https://app.zenhub.com/workspaces/tfgrag-68d94a186434450034791a44/board}{\eng{ZenHub}} mediante la utilización de un tablero \eng{kanban}, los \eng{story points} y de los \eng{Milestone Burdown}. Este proyecto asociado al de \eng{GitHub}.
		
		El progreso del proyecto también puede verse en un~\myhref{https://github.com/users/lbr1014/projects/1}{\eng{GitHub Project}} asociado al repositorio de \eng{GitHub}, en el cual puede observarse una tabla con las \eng{issues} de cada \eng{sprint} y un tablero \eng{kanban}. Además en el apartado~\myhref{https://github.com/users/lbr1014/projects/1/insights}{\eng{insights}} se han generado los \eng{Burndown chart} de cada \eng{sprint}.
		\item Informe de calidad: para analizar la calidad del código de la aplicación definitiva (la que se encuentra dentro de la carpeta \eng{PythIA} en \eng{GitHub}) se ha utilizado~\myhref{https://sonarcloud.io/project/overview?id=lbr1014_TFG_RAG}{\eng{SonarCloud}}.
\end{itemize}